\documentclass[11pt]{article}
\usepackage[margin=1.1in]{geometry}
\usepackage{amsmath, amssymb}
\usepackage{booktabs}
\usepackage{graphicx}
\usepackage[hidelinks]{hyperref}
\usepackage{caption}
\captionsetup{font=small, labelfont=bf}

\newcommand{\rhat}{\widehat{RC}}
\newcommand{\that}[1]{\widehat{\theta}_{#1}}
\newcommand{\mstar}{m^{*}}

\title{Endogenous Attention and the Hierarchy of Invariance\\
\large When As-If Estimation Breaks, What Audits Cannot See,\\
and Why Depth Must Be Bought Wholesale}
\author{Briac Sockalingum\thanks{briac@berkeley.edu. First draft --- to be
completed. Companion to ``Do Structural Models Recover Mechanisms?''
(2026), \url{https://github.com/briacSck/structural-interp}. Code and
results: \url{https://github.com/briacSck/endogenous-attention}. Hypotheses
were pre-registered by commit before execution
(\texttt{design/research\_design.md}).}}
\date{July 2026 --- Working note (first draft)}

\begin{document}
\maketitle

\begin{abstract}
\noindent In a companion paper we showed that structural estimation of
dynamic discrete choice is robust to bounded rationality that stays inside
the model class: agents with \emph{fixed} inattention are absorbed into
as-if parameters, and counterfactuals survive. This note breaks the
assumption doing the work. When attention is \emph{endogenous} --- chosen
against stakes, as rational-inattention theory requires --- a policy that
changes stakes moves the cognitive primitive itself. In a laboratory of
neural sparse-max agents whose attention cost $\kappa$ is known by
construction, we find, as pre-registered: (i) as-if counterfactual failure
tracks the policy-induced attention repricing $|m^*_1 - m^*_0|$ (Spearman
0.88), vanishing at both corners of the attention-cost grid; (ii) internal
conformity audits computed pre-policy are \emph{blind} to the coming
failure (Spearman $-0.09$) --- correctly so, since every agent's internals
genuinely conform before the intervention; the failure lives one level up.
(iii) The deep primitive is identifiable: minimum distance on the
cross-market profile of as-if attention slopes recovers the attention cost
exactly ($\widehat{\kappa} = 0.200$, truth $0.200$). But our pre-registered
hypothesis that plug-in deep corrections beat as-if predictions is
\emph{rejected}: swapping a deeper $\theta_2$ into a jointly-estimated
as-if vector breaks the compensation among its coordinates
(the pseudo-true-value logic), and predicts worse in four of five markets.
The hierarchy of invariance --- behavioral parameters $<$ as-if preferences
$<$ cognitive primitives --- thus comes with a warning label: each policy
class requires estimating one level deeper, and depth must be bought
wholesale, by joint estimation with the attention technology inside the
likelihood, not retail, by parameter swapping. The coherent joint estimator
is this paper's designated completion.
\end{abstract}

\section{Introduction}\label{sec:intro}

The companion paper's reassuring finding --- inattention is absorbed into
as-if parameters and counterfactuals survive, an empirical realization of
the \citet{fosgerau2020} equivalence --- rested on one assumption: the
cognitive primitive stays fixed under the intervention. Rational-inattention
theory says it should not. Attention responds to stakes
\citep{matejkamckay2015, gabaix2014, mackowiak2023}; a subsidy that changes
payoff differences changes what is worth attending to. If so, as-if
parameters estimated pre-policy are not invariant to the policy --- the
Lucas critique \citep{lucas1976}, reconstituted one level up, for cognitive
rather than behavioral parameters.

This note measures that reconstitution in a laboratory where every level of
the hierarchy is observable. Our agents choose attention $\mstar(\kappa, s)$
by trading the exact value of perceiving a cost-regime state against a
quadratic attention cost $\kappa$; stakes $s$ vary across markets; the
truth --- $(\theta_1, \theta_2, RC, \kappa)$ and the induced attention ---
is known by construction. Three pre-registered questions: does as-if
counterfactual failure track the policy-induced attention repricing (H1)?
Can pre-policy internal audits --- the companion paper's conformity
diagnostics --- see it coming (H2)? And does recovering the deep primitive
restore counterfactual validity (H3)?

The answers are yes, no, and --- instructively --- \emph{not the way we
pre-registered}. H1: failure tracks repricing at Spearman 0.88, with the
pre-registered corner structure (at $\kappa \to 0$ and $\kappa \to \infty$
attention is pinned, nothing reprices, as-if predictions are fine). H2:
the audits are blind (Spearman $-0.09$), and \emph{correctly} blind ---
pre-policy internals genuinely conform to the model class; what breaks is
the primitive's response to the intervention, which no pre-policy internal
measurement contains. This measures precisely what mechanism audits
certify: invariance under mechanism-preserving interventions, nothing
stronger. H3: the attention cost is exactly identifiable from
cross-sectional stakes variation ($\widehat{\kappa} = 0.200$, truth 0.200)
--- but plug-in deep corrections \emph{underperform} the as-if benchmark in
four of five markets, in both variants we tried. The diagnosis is the
pseudo-true-value logic of quasi-maximum likelihood \citep{white1982}: the
as-if vector's coordinates are jointly chosen and their biases compensate;
replacing one coordinate with a better estimate of its ``true'' value
breaks the compensation. Depth is necessary (H1) and identifiable (H3a),
but must be bought wholesale --- joint estimation of
$(\theta_1, \theta_2, RC, \kappa)$ with the attention technology inside the
likelihood --- not retail. Section~\ref{sec:joint} specifies that estimator
as the paper's designated completion.

\section{The laboratory}\label{sec:lab}

\paragraph{Environment.} The companion paper's two-state-variable bus
replacement problem: mileage $x$ (60 bins) and a persistent observable cost
regime $c$ (5 states) scaling maintenance costs by $\theta_2$. Calibration
for this paper (chosen by a pre-registered instrument-design step
\emph{before} any failure metric was computed): $\theta_2 = 1.5$, regime
persistence $0.95$, $(\theta_1, RC, \beta) = (0.05, 10, 0.95)$. The policy
intervention is a 50\% proportional replacement subsidy ($RC\!: 10 \to 5$).

\paragraph{Endogenous attention.} Agents perceive the regime as
$c' = m c + (1 - m)\bar c$ and act optimally on the perceived state; the
attention level is chosen as
\[
\mstar(\kappa, \text{env}) \;=\; \arg\max_{m \in [0,1]}\;
B(m; \text{env}) - \kappa m^2,
\]
where $B(m)$ --- the expected discounted payoff of the $m$-perception
policy --- is computed exactly by policy evaluation on the solved model.
Comparative statics are verified by unit test: $\mstar$ increases in
stakes, decreases in $\kappa$, with corners $\mstar(0) = 1$ and
$\mstar(\infty) = 0$. Under the calibration, the subsidy reprices attention
by up to $|\mstar_1 - \mstar_0| = 0.175$ across the $\kappa$ grid, with
repricing exactly zero at both corners --- the instrument H1 requires.

\paragraph{Agents, estimation, audits.} For each $\kappa$ on a ten-point
grid, an MLP agent (companion architecture and seeds) is trained on the
$\mstar(\kappa)$-perception policy; the econometrician estimates
$(\theta_1, \theta_2, RC)$ by NFXP on its simulated pre-policy panel;
adaptation to the subsidy follows the \emph{stable-cognitive-primitive}
rule ($\kappa$ fixed, attention re-priced to $\mstar_1$, policy re-derived
and re-learned); the as-if prediction solves the estimated model at
$(\that{1}, \that{2}, \rhat/2)$. The conformity audits (per-channel
dispersion of locally patched causal sensitivities with the insensitivity
gate) are ported unchanged from the companion paper and computed on
pre-policy activations only.

\section{Results}\label{sec:results}

\subsection{H1 (confirmed): failure tracks attention repricing}
Across the $\kappa$ grid, as-if counterfactual policy RMSE correlates with
the attention repricing $|\mstar_1 - \mstar_0|$ at Spearman
\textbf{0.88}. The corner structure is exact: for $\kappa \le 0.05$,
attention is pinned at $m = 1$ pre- and post-subsidy and the as-if
prediction is at the clone baseline (RMSE $\approx 0.0045$); in the
interior, failure peaks at 3.6$\times$ baseline ($\kappa = 0.1$:
$|\Delta m| = 0.15$, RMSE 0.0161). The estimated $\that{2}$ tracks
$\mstar_0$ across the grid (1.53 at $m = 1$ down to 0.19 at $m = 0.2$):
the companion paper's as-if absorption, reproduced --- and now exhibited as
exactly the thing that breaks when the primitive moves.

\emph{Honest magnitudes.} These repricing failures are real but moderate
--- an order of magnitude smaller than the companion paper's out-of-class
categorical failures (max RMSE 0.016 vs.\ 0.25; replacement-rate errors
$\sim$0.001). Under this calibration, moving the primitive hurts less than
misspecifying the mechanism class. The contribution is the pattern and its
logic, not the raw bias size.

\subsection{H2 (confirmed): pre-policy audits are blind, and rightly so}
The conformity audit is uncorrelated with the coming failure (Spearman
$\mathbf{-0.09}$; audit scores flat across the grid, $cv_x \in
[0.69, 0.79]$, $cv_c \in [0.53, 0.64]$). This is not an audit defect: every
fixed-$m$ sparse-max agent's internal mechanism genuinely conforms to the
assumed class pre-policy, and the audit correctly says so. The failure is
generated by the primitive's response to the intervention --- information
that does not exist in pre-policy internals. Together with the companion
paper this gives the audit technology a precise warranty: \emph{conformity
audits certify counterfactual validity under mechanism-preserving
interventions; under primitive-moving interventions they are silent, not
wrong.}

\subsection{H3 (identification confirmed; plug-in correction rejected)}
\label{sec:h3}
A cross-section of five markets with stakes shifters
$s \in \{0.5, \dots, 2.0\}$ (payoff scale, shock scale fixed) makes the
as-if attention slope observable: $\that{2}(s)$ rises monotonically with
stakes (0.54 to 1.18) because higher-stakes agents attend more. Minimum
distance on this profile, using the known attention technology, recovers
the deep primitive \emph{exactly}: $\widehat{\kappa} = 0.200$ against a
truth of $0.200$. The hierarchy's second level is identifiable from
standard cross-sectional variation.

But the pre-registered H3 --- deep counterfactuals beat as-if by an order
of magnitude --- is \textbf{rejected}, and we report it as pre-registered.
Plug-in deep predictions underperform as-if in four of five markets, in
both variants (pooled $\theta_2^{deep}$, and per-market de-attenuation
$\that{2}(s)/\widehat{m}_0$). The diagnosis: the as-if vector is a
pseudo-true value \citep{white1982} --- its coordinates are chosen jointly
by maximum likelihood, so their individual biases compensate in prediction.
Swapping one coordinate for a deeper estimate, however well identified,
produces a vector coherent at no level, which predicts worse than honest
misspecification. Depth must be bought wholesale, not retail.

\section{The designated completion: a coherent deep estimator}
\label{sec:joint}
The constructive resolution this paper still owes --- and will include, in
this paper rather than a sequel --- is joint full-solution estimation of
$(\theta_1, \theta_2, RC, \kappa)$: for each candidate parameter vector,
solve the model, solve the attention problem $\mstar(\kappa, s)$ per
market, form the implied perception policies, and evaluate the pooled
likelihood across markets. This estimator is internally coherent at the
deep level by construction; the computational burden (an attention fixed
point inside the likelihood) is the reason it is a designated next
experiment rather than a footnote. Its pre-registration will follow the
protocol used throughout: hypotheses committed before execution.

\section{Discussion}\label{sec:discussion}
Three sentences of positioning. For structural practice: the Lucas critique
does not stop at preferences; whenever a policy moves stakes, attention-like
primitives move, and the practitioner's choice is between as-if estimates
with a known invariance warranty and deep estimates bought coherently. For
interpretability: audit warranties are bounded by mechanism invariance ---
a boundary this laboratory measures rather than conjectures. For the
program: the hierarchy of invariance gives the ``microeconomics of
artificial agents'' its organizing principle --- ask, of any estimator and
any audit, \emph{which level it certifies}.

Limitations mirror the companion paper (toy scale, MLP agents, transitions
known), plus two specific to this note: the attention technology is known
to the analyst in E3 (identification uses its functional form), and the
quadratic-cost sparse-max is one of several attention models --- the exact
RI-logit variant is a designated robustness exercise.

\begin{thebibliography}{99}\small

\bibitem[Fosgerau et al.(2020)]{fosgerau2020}
Fosgerau, M., Melo, E., de Palma, A., \& Shum, M. (2020). Discrete choice
and rational inattention: A general equivalence result.
\emph{International Economic Review}, 61(4), 1569--1589.

\bibitem[Gabaix(2014)]{gabaix2014}
Gabaix, X. (2014). A sparsity-based model of bounded rationality.
\emph{Quarterly Journal of Economics}, 129(4), 1661--1710.

\bibitem[Lucas(1976)]{lucas1976}
Lucas, R.~E. (1976). Econometric policy evaluation: A critique.
\emph{Carnegie-Rochester Conference Series on Public Policy}, 1, 19--46.

\bibitem[Ma\'ckowiak et al.(2023)]{mackowiak2023}
Ma\'ckowiak, B., Mat\v{e}jka, F., \& Wiederholt, M. (2023). Rational
inattention: A review. \emph{Journal of Economic Literature}, 61(1),
226--273.

\bibitem[Mat\v{e}jka \& McKay(2015)]{matejkamckay2015}
Mat\v{e}jka, F., \& McKay, A. (2015). Rational inattention to discrete
choices: A new foundation for the multinomial logit model. \emph{American
Economic Review}, 105(1), 272--298.

\bibitem[Sockalingum(2026)]{companion2026}
Sockalingum, B. (2026). Do structural models recover mechanisms? Evidence
from neural agents with observable internals. Working note,
\url{https://github.com/briacSck/structural-interp}.

\bibitem[White(1982)]{white1982}
White, H. (1982). Maximum likelihood estimation of misspecified models.
\emph{Econometrica}, 50(1), 1--25.

\end{thebibliography}

% TODO (author + feynman pass): verify White (1982) page range; add
% Caplin, Dean & Leahy (JPE 2022?) for attention-cost identification if it
% survives verification; inference (permutation p / CIs) on the 0.88 once
% the kappa grid is final; figure e1_e2.png inclusion.

\end{document}
